\begin{textsample}{NEG Dim 1 – pb – Score 3.00 – pb/pb\_ef\_20.txt}  \label{ex:f1_neg_169}
5.2.6.4 Avaliação A avaliação no Ensino de Teatro, assim como a de outras linguagens artísticas, deverá ser de cunho teórico-prático, observando as especificidades dessa linguagem.
Segundo os PCN de Arte, o(a) professor(a) de Teatro deverá ter como critério para a avaliação os pontos a seguir : - Saber improvisar e atuar nas situações de jogos, explorando as capacidades do corpo e da voz. - Estar capacitado para criar cenas escritas ou encenadas, reconhecendo e organizando os recursos para a sua estruturação. - Estar capacitado a emitir opiniões sobre a atividade teatral, com clareza e critérios fundamentados, sem discriminação estética, artística, étnica ou de gênero. - Identificar momentos importantes da história do teatro, da dramaturgia local, nacional ou internacional, refletindo e relacionando os aspectos estéticos e cênicos. - Valorizar as fontes de documentação, os acervos e os arquivos da produção artística teatral. ( p. 99, Parâmetros Curriculares Nacionais - ARTE ) Tomando como base os pontos elencados acima, o(a) professor(a) deve ter claro que avaliar é uma ação processual, além de considerar o desenvolvimento do(a) estudante em sua individualidade na sua trajetória.
Não deve esperar, portanto, que o processo avaliativo gere apenas notas quantitativas que levem à progressão do(a) estudante, mas sim centrar o fazer pedagógico nele, para que, verdadeiramente, a aprendizagem se concretize.
O(A) professor(a) deverá considerar todo o percurso do(a) estudante e sua construção ao longo desse processo.

% matched lemmas: 
\end{textsample}
